% =============================================================================
%  Habivi — Diseño de Base de Datos (Hive)
%  Diapositiva horizontal para Overleaf (beamer, landscape)
% =============================================================================
\documentclass[11pt,landscape]{beamer}
\usepackage[utf8]{inputenc}
\usepackage[T1]{fontenc}
\usepackage[spanish]{babel}
\usepackage{booktabs}
\usepackage{array}
\usepackage{xcolor}
\usepackage{colortbl}
\usepackage{tikz}
\usetikzlibrary{shapes,arrows,positioning,fit}

% ---- Tema minimalista oscuro ----
\setbeamercolor{background canvas}{bg=black!90}
\setbeamercolor{normal text}{fg=white}
\setbeamercolor{frametitle}{fg=white}
\setbeamercolor{title}{fg=white}
\setbeamercolor{structure}{fg=cyan!70}
\setbeamercolor{alerted text}{fg=orange}
\setbeamercolor{footline}{fg=gray}

\setbeamertemplate{navigation symbols}{}
\setbeamertemplate{frametitle}{\vspace{0.3cm}\insertframetitle\par}

\title{\textbf{\Huge Habivi}}
\subtitle{\large Diseño de Base de Datos Local --- Hive}
\author{Arquitectura de Software}
\date{}

\begin{document}

% ---- Portada ----
\begin{frame}
  \titlepage
  \vfill
  \centering
  {\small \color{gray} App de productividad personal con gamificación}
\end{frame}

% ---- Slide 1: Visión general ----
\begin{frame}{Visión General}
  \textbf{¿Por qué Hive?}
  \begin{itemize}
    \item Base de datos NoSQL clave--valor, \alert{100\,\% local} (sin servidor)
    \item Sin reflectividad ni código nativo pesado (ideal para Flutter)
    \item Escritura ultra rápida en memoria + persistencia asíncrona
    \item \texttt{localStorage} en web; archivos binarios en móvil/desktop
  \end{itemize}

  \vspace{0.4cm}
  \textbf{Principios de diseño}
  \begin{enumerate}
    \item Una \alert{caja} por entidad, 12 cajas en total
    \item \texttt{typeId} fijo (0–11) para compatibilidad hacia atrás
    \item \texttt{flush()} tras cada escritura (create, update, delete)
    \item Cierre de Hive solo en \texttt{AppLifecycleState.detached}
  \end{enumerate}
\end{frame}

% ---- Slide 2: Tabla resumen ----
\begin{frame}{Entidades y Cajas}
  \centering
  \rowcolors{2}{black!75}{black!60}
  \begin{tabular}{ccllc}
    \toprule
    \rowcolor{black!45}
    \textbf{typeId} & \textbf{Caja} & \textbf{Modelo} & \textbf{Adapter} & \textbf{Campos} \\
    \midrule
    0  & usuarioBox      & Usuario        & auto       & 8 \\
    1  & cuentaBox       & Cuenta         & auto       & 5 \\
    2  & habitoBox       & Habito         & auto       & 9 \\
    3  & registroHabitoBox & RegistroHabito & auto     & 4 \\
    4  & tareaBox        & Tarea          & auto       & 7 \\
    5  & sesionPomodoroBox & SesionPomodoro & auto     & 6 \\
    6  & notaBox         & Nota           & auto       & 2 \\
    7  & eventoBox       & Evento         & auto       & 3 \\
    8  & backupBox       & Backup         & auto       & 3 \\
    9  & sesionEstudioBox & SesionEstudio  & \alert{manual} & 6 \\
    10 & logroBox        & Logro          & \alert{manual} & 2 \\
    11 & rachaBox        & Racha          & \alert{manual} & 7 \\
    \bottomrule
  \end{tabular}
\end{frame}

% ---- Slide 3: Diagrama entidad-relación conceptual ----
\begin{frame}{Diagrama de Entidades}
  \centering
  \begin{tikzpicture}[
    box/.style={rectangle, rounded corners, draw=cyan!60, fill=cyan!15,
                minimum width=4.2cm, minimum height=0.7cm, font=\small},
    rel/.style={-latex, draw=orange!70, thick},
    every node/.style={font=\footnotesize}
  ]

    % Fila superior: Usuario + Cuenta
    \node[box] (usr) {Usuario (idUsuario*)};
    \node[box, right=1.2cm of usr] (cta) {Cuenta (idUsuario*)};

    % Fila 2: Hábitos y tareas
    \node[box, below=0.9cm of usr] (hab) {Hábito};
    \node[box, right=1.2cm of hab] (reg) {RegistroHábito};
    \node[box, below=0.9cm of hab] (tar) {Tarea};

    % Fila 3: Sesiones
    \node[box, right=1.2cm of tar] (pom) {SesiónPomodoro};
    \node[box, below=0.9cm of tar] (est) {SesiónEstudio};

    % Fila 4: Notas, Eventos, Backup
    \node[box, right=1.2cm of est] (nota) {Nota};
    \node[box, below=0.9cm of est] (evt) {Evento};
    \node[box, right=1.2cm of evt] (bkp) {Backup};

    % Fila 5: Logros y Rachas
    \node[box, below=0.9cm of evt] (log) {Logro};
    \node[box, right=1.2cm of log] (rch) {Racha};

    % Relaciones conceptuales
    \draw[rel] (usr) -- (cta);
    \draw[rel, dashed] (usr) |- (hab);
    \draw[rel, dashed] (hab) -- (reg);
    \draw[rel, dashed] (usr) |- (tar);
    \draw[rel, dashed] (usr) |- (est);
    \draw[rel, dashed] (usr) |- (log);
    \draw[rel, dashed] (usr) |- (rch);

    % Nota al pie
    \node[below=0.2cm of rch, xshift=-2.5cm, text width=10cm, align=center,
          font=\tiny\color{gray}] {
      * Relaciones conceptuales — Hive no impone claves foráneas. \\
      Las relaciones se resuelven en código (repositorios y servicios).
    };
  \end{tikzpicture}
\end{frame}

% ---- Slide 4: Detalle de campos ----
\begin{frame}{Campos por Entidad (I)}
  \centering
  \rowcolors{2}{black!75}{black!60}
  \begin{tabular}{clll}
    \toprule
    \rowcolor{black!45}
    \textbf{Modelo} & \textbf{Campo} & \textbf{Tipo} & \textbf{Descripción} \\
    \midrule
    \multirow{8}{*}{Usuario} & idUsuario & int & ID único \\
     & nombre & String & \\
     & apellido & String & \\
     & energia & int & Energía actual (0--100) \\
     & energiaMax & int & Máxima energía \\
     & puntosProductividad & int & Puntos acumulados \\
     & estadoPersonaje & String & Estado del personaje \\
     & fechaInicio & String & Fecha de registro \\
    \midrule
    \multirow{5}{*}{Cuenta} & idUsuario & int & FK conceptual \\
     & nickname & String & \\
     & contrasena & String & \\
     & correo & String & \\
     & telefono & int & \\
    \midrule
    \multirow{9}{*}{Hábito} & idHabito & int & \\
     & nombreHabito & String & \\
     & atributo & String & \\
     & tipo & String & ``aspecto|iconoKey'' \\
     & completadoHoy & bool & \\
     & fechaUltimoCompletado & String & \\
     & fechasCompletadas & List<String> & Historial \\
     & vecesPorSemana & int? & Frecuencia objetivo \\
     & fechaCreacion & String? & \\
    \bottomrule
  \end{tabular}
\end{frame}

\begin{frame}{Campos por Entidad (y II)}
  \centering
  \rowcolors{2}{black!75}{black!60}
  \begin{tabular}{clll}
    \toprule
    \rowcolor{black!45}
    \textbf{Modelo} & \textbf{Campo} & \textbf{Tipo} & \textbf{Descripción} \\
    \midrule
    \multirow{4}{*}{RegistroHábito} & idRegistro & int & \\
     & fecha & String & \\
     & completado & bool & \\
     & impacto & int & \\
    \midrule
    \multirow{7}{*}{Tarea} & idTarea & int & \\
     & nombreTarea & String & \\
     & puntaje & int & \\
     & metadata & String & Legado ``|fecha|notas'' \\
     & fecha & String & \\
     & notas & String & \\
     & completada & bool & \\
    \midrule
    \multirow{6}{*}{SesiónPomodoro} & idSesion & int & \\
     & fechaInicio & String & \\
     & fechaFinal & String & \\
     & duracion & int & Minutos \\
     & completada & bool & \\
     & tipoSesion & String & pomodoro / descanso \\
    \midrule
    \multirow{2}{*}{Nota} & idNota & int & \\
     & valor & String & Contenido \\
    \bottomrule
  \end{tabular}
\end{frame}

\begin{frame}{Campos por Entidad (y III)}
  \centering
  \rowcolors{2}{black!75}{black!60}
  \begin{tabular}{clll}
    \toprule
    \rowcolor{black!45}
    \textbf{Modelo} & \textbf{Campo} & \textbf{Tipo} & \textbf{Descripción} \\
    \midrule
    \multirow{3}{*}{Evento} & idEvento & int & \\
     & fechaEvento & String & \\
     & finalizado & bool & \\
    \midrule
    \multirow{3}{*}{Backup} & idBackup & int & \\
     & fechaBackup & String & \\
     & descripcion & String & \\
    \midrule
    \multirow{6}{*}{SesiónEstudio} & idSesion & int & \\
     & metodo & String & pomodoro, feynman, etc. \\
     & duracionMinutos & int & \\
     & puntosObtenidos & int & \\
     & fecha & String & \\
     & nota & String & \\
    \midrule
    \multirow{2}{*}{Logro} & idLogro & String & Clave del logro \\
     & fechaDesbloqueo & String & \\
    \midrule
    \multirow{7}{*}{Racha} & idRacha & int & \\
     & tipo & String & diaria / semanal \\
     & cantidad & int & Días o semanas \\
     & fechaUltimoCompletado & String & \\
     & enRiesgo & bool & Solo semanal \\
     & fechaInicioPeriodoRecuperacion & String & \\
     & fechaCreacion & String? & \\
    \bottomrule
  \end{tabular}
\end{frame}

% ---- Slide 5: Estrategia de persistencia ----
\begin{frame}{Estrategia de Persistencia}
  \textbf{Escritura inmediata}
  \begin{itemize}
    \item Cada repositorio llama a \texttt{box.flush()} tras \texttt{add},
          \texttt{put}, \texttt{delete} y \texttt{clear}
    \item Garantiza que los datos sobrevivan incluso si la app se cierra
          abruptamente
  \end{itemize}

  \vspace{0.3cm}
  \textbf{Ciclo de vida}
  \begin{itemize}
    \item \texttt{main.dart}: \texttt{AppDatabase.initialize()} abre las 12 cajas
    \item \texttt{\_AppLifecycleWrapper}: cierra Hive solo en
          \texttt{AppLifecycleState.detached} (no en \texttt{paused})
    \item En web: se fuerza puerto \texttt{5000} para que
          \texttt{localStorage} (por origen) sea consistente entre ejecuciones
  \end{itemize}

  \vspace{0.3cm}
  \textbf{Manejo de versiones}
  \begin{itemize}
    \item \texttt{typeId} fijo: cada modelo tiene un ID de adaptador
          permanente (0--11)
    \item Nuevos campos: se añaden al final con el siguiente \texttt{HiveField}
          (ej. Habito saltó de 7 a 9 campos)
    \item Adapters manuales (SesionEstudio, Logro, Racha): control total sobre
          serialización, sin dependencia de generación de código
  \end{itemize}
\end{frame}

% ---- Slide final ----
\begin{frame}{}

  \vfill
  \centering
  {\Huge \textbf{Gracias}}
  \vspace{0.5cm}

  {\large \url{https://github.com/jladinori/Habivi}}
  \vfill

\end{frame}

\end{document}
